\documentclass[11pt,a4paper]{article}

\usepackage[utf8]{inputenc}
\usepackage[T1]{fontenc}
\usepackage[french]{babel}
\usepackage[a4paper,margin=2.3cm]{geometry}
\usepackage{amsmath,amssymb,amsthm,mathtools}
\usepackage{lmodern}
\usepackage{enumitem}
\usepackage{hyperref}
\usepackage{xcolor}
\usepackage{fancyhdr}
\usepackage{stmaryrd}
\usepackage{mathrsfs}
\usepackage{array}

\hypersetup{
    colorlinks=true,
    linkcolor=black,
    urlcolor=blue,
    citecolor=black
}

\setlength{\parindent}{0pt}
\setlength{\parskip}{0.6em}

\pagestyle{fancy}
\fancyhf{}
\fancyhead[L]{Algèbre linéaire --- Chapitre 3}
\fancyhead[R]{Systèmes linéaires}
\fancyfoot[C]{\thepage}

% Environnements
\newtheorem{definition}{Définition}[section]
\newtheorem{theoreme}[definition]{Théorème}
\newtheorem{proposition}[definition]{Proposition}
\newtheorem{corollaire}[definition]{Corollaire}
\newtheorem{lemme}[definition]{Lemme}
\theoremstyle{remark}
\newtheorem{remarque}[definition]{Remarque}
\newtheorem{exemple}[definition]{Exemple}
\newtheorem*{methode}{Méthode}
\newtheorem*{idee}{Idée}

% Commandes
\newcommand{\R}{\mathbb{R}}
\newcommand{\C}{\mathbb{C}}
\newcommand{\K}{\mathbb{K}}
\newcommand{\Vect}{\mathrm{Vect}}
\newcommand{\rg}{\mathrm{rg}}

\begin{document}

\begin{center}
    {\LARGE \textbf{Chapitre 3 --- Systèmes linéaires}}\\[0.5em]
    {\large Résoudre, interpréter, structurer}
\end{center}

\hrule
\vspace{1em}

\tableofcontents

\newpage

\section{Pourquoi ce chapitre est-il fondamental ?}

Les systèmes linéaires sont l’un des points de départ naturels de l’algèbre linéaire. Ils apparaissent partout :
\begin{itemize}
    \item en géométrie, pour décrire des intersections de droites ou de plans ;
    \item en calcul matriciel ;
    \item dans l’étude des espaces vectoriels et des applications linéaires ;
    \item en sciences physiques, en économie, en informatique, dès qu’on modélise des relations linéaires entre inconnues.
\end{itemize}

\begin{idee}
Résoudre un système linéaire, c’est chercher les valeurs des inconnues qui satisfont simultanément plusieurs équations linéaires.
\end{idee}

\begin{remarque}
Ce chapitre est central, car beaucoup de questions d’algèbre linéaire se ramènent, en pratique, à la résolution d’un système linéaire :
\begin{itemize}
    \item tester si une famille est libre ;
    \item déterminer si une famille est génératrice ;
    \item calculer des coordonnées ;
    \item résoudre des problèmes d’intersection de sous-espaces.
\end{itemize}
\end{remarque}

\section{Définition d’un système linéaire}

\subsection{Équation linéaire}

\begin{definition}
Soient \(x_1,\dots,x_n\) des inconnues à valeurs dans un corps \(\K\) (le plus souvent \(\R\) ou \(\C\)).  
Une \textbf{équation linéaire} en \(x_1,\dots,x_n\) est une équation de la forme
\[
a_1x_1+\cdots+a_nx_n=b,
\]
où \(a_1,\dots,a_n,b\in\K\).
\end{definition}

\begin{remarque}
Le mot \og linéaire \fg{} signifie que les inconnues apparaissent seulement au premier degré, sans produit entre elles, sans puissance, sans fonction non linéaire.
\end{remarque}

\begin{exemple}
\[
2x-3y+z=5
\]
est une équation linéaire en \(x,y,z\).

En revanche,
\[
x^2+y=1
\qquad \text{ou} \qquad
xy+z=0
\]
ne sont pas des équations linéaires.
\end{exemple}

\subsection{Système linéaire}

\begin{definition}
Un \textbf{système linéaire} de \(p\) équations à \(n\) inconnues est un ensemble d’équations linéaires du type
\[
\begin{cases}
a_{11}x_1+\cdots+a_{1n}x_n=b_1,\\
a_{21}x_1+\cdots+a_{2n}x_n=b_2,\\
\hspace{2.5cm}\vdots\\
a_{p1}x_1+\cdots+a_{pn}x_n=b_p.
\end{cases}
\]
\end{definition}

\begin{definition}
Une \textbf{solution} du système est un \(n\)-uplet \((x_1,\dots,x_n)\in\K^n\) qui vérifie simultanément toutes les équations du système.
\end{definition}

\begin{exemple}
Le système
\[
\begin{cases}
x+y=3,\\
2x-y=0
\end{cases}
\]
admet pour solution \((1,2)\), car
\[
1+2=3
\qquad \text{et} \qquad
2\cdot 1-2=0.
\]
\end{exemple}

\section{Systèmes homogènes et non homogènes}

\subsection{Définitions}

\begin{definition}
Un système linéaire est dit \textbf{homogène} si tous les seconds membres sont nuls :
\[
\begin{cases}
a_{11}x_1+\cdots+a_{1n}x_n=0,\\
\hspace{2.2cm}\vdots\\
a_{p1}x_1+\cdots+a_{pn}x_n=0.
\end{cases}
\]
\end{definition}

\begin{definition}
Un système linéaire est dit \textbf{non homogène} s’il possède au moins un second membre non nul.
\end{definition}

\begin{exemple}
Le système
\[
\begin{cases}
x+2y-z=0,\\
3x-y+4z=0
\end{cases}
\]
est homogène.

Le système
\[
\begin{cases}
x+2y-z=1,\\
3x-y+4z=0
\end{cases}
\]
ne l’est pas.
\end{exemple}

\subsection{Propriété fondamentale des systèmes homogènes}

\begin{proposition}
Tout système linéaire homogène admet au moins une solution : la solution nulle
\[
(x_1,\dots,x_n)=(0,\dots,0).
\]
\end{proposition}

\begin{proof}
Si toutes les inconnues valent \(0\), alors chaque membre de gauche est nul, donc chaque équation est vérifiée.
\end{proof}

\begin{remarque}
C’est une différence importante avec les systèmes non homogènes, qui peuvent ne pas avoir de solution du tout.
\end{remarque}

\section{Interprétation géométrique}

\subsection{Dans le plan}

\begin{exemple}
Dans \(\R^2\), une équation linéaire
\[
ax+by=c
\]
représente une droite affine lorsque \((a,b)\neq (0,0)\).

Ainsi, résoudre un système de deux équations à deux inconnues revient à chercher l’intersection de deux droites.
\end{exemple}

\begin{remarque}
Trois situations peuvent se produire :
\begin{itemize}
    \item une unique solution : les droites sont sécantes ;
    \item aucune solution : les droites sont parallèles distinctes ;
    \item une infinité de solutions : les droites sont confondues.
\end{itemize}
\end{remarque}

\subsection{Dans l’espace}

\begin{exemple}
Dans \(\R^3\), une équation linéaire
\[
ax+by+cz=d
\]
représente un plan affine lorsque \((a,b,c)\neq (0,0,0)\).

Résoudre un système de deux ou trois équations revient alors à étudier l’intersection de plans.
\end{exemple}

\section{Opérations élémentaires sur les équations}

\subsection{Principe}

\begin{idee}
Pour résoudre un système, on peut le transformer en un système plus simple, à condition de ne pas changer l’ensemble des solutions.
\end{idee}

\begin{definition}
On appelle \textbf{opérations élémentaires sur les lignes} les opérations suivantes :
\begin{enumerate}[label=\arabic*)]
    \item échanger deux équations ;
    \item multiplier une équation par un scalaire non nul ;
    \item ajouter à une équation un multiple d’une autre.
\end{enumerate}
\end{definition}

\begin{proposition}
Les opérations élémentaires conservent l’ensemble des solutions d’un système.
\end{proposition}

\begin{proof}
\textbf{1.} Échanger deux équations ne change évidemment pas les conditions imposées aux inconnues.

\textbf{2.} Multiplier une équation par un scalaire non nul ne change pas l’ensemble des solutions, car une égalité
\[
L=b
\]
équivaut à
\[
\lambda L=\lambda b
\]
si \(\lambda\neq 0\).

\textbf{3.} Remplacer une équation \(L_2=b_2\) par
\[
L_2+\lambda L_1=b_2+\lambda b_1
\]
ne change pas l’ensemble des solutions, car toute solution de l’ancien système vérifie la nouvelle équation, et réciproquement, connaissant \(L_1=b_1\), on retrouve \(L_2=b_2\) en soustrayant \(\lambda L_1\).
\end{proof}

\begin{remarque}
C’est ce principe qui justifie toute la méthode du pivot de Gauss.
\end{remarque}

\section{Méthode du pivot de Gauss}

\subsection{Objectif}

\begin{idee}
La méthode du pivot de Gauss consiste à transformer le système en un système équivalent plus simple, de forme triangulaire, afin de résoudre successivement les inconnues.
\end{idee}

\subsection{Exemple simple}

\begin{exemple}
Résolvons le système
\[
\begin{cases}
x+y=3,\\
2x-y=0.
\end{cases}
\]

On conserve la première ligne, et on remplace la deuxième par
\[
L_2 \leftarrow L_2-2L_1.
\]
On obtient
\[
\begin{cases}
x+y=3,\\
-3y=-6.
\end{cases}
\]
Donc
\[
y=2.
\]
En reportant dans la première équation :
\[
x+2=3 \Longrightarrow x=1.
\]
La solution unique est
\[
(x,y)=(1,2).
\]
\end{exemple}

\subsection{Forme triangulaire}

\begin{definition}
Un système est dit \textbf{triangulaire} si les inconnues apparaissent progressivement, par exemple :
\[
\begin{cases}
a_{11}x_1+a_{12}x_2+\cdots+a_{1n}x_n=b_1,\\
a_{22}x_2+\cdots+a_{2n}x_n=b_2,\\
\hspace{2cm}\vdots\\
a_{nn}x_n=b_n.
\end{cases}
\]
\end{definition}

\begin{remarque}
Un système triangulaire se résout par \textbf{remontée} :
\begin{itemize}
    \item on détermine d’abord la dernière inconnue ;
    \item puis on remonte ligne par ligne.
\end{itemize}
\end{remarque}

\subsection{Description générale de la méthode}

\begin{methode}
Pour résoudre un système linéaire par pivot de Gauss :
\begin{enumerate}
    \item on choisit une inconnue à éliminer ;
    \item on utilise une équation dans laquelle son coefficient est non nul comme \textbf{ligne pivot} ;
    \item on élimine cette inconnue dans les autres équations ;
    \item on recommence sur le système réduit ;
    \item une fois la forme triangulaire obtenue, on remonte.
\end{enumerate}
\end{methode}

\section{Systèmes compatibles, incompatibles, déterminés, indéterminés}

\subsection{Définitions}

\begin{definition}
Un système est dit \textbf{compatible} s’il admet au moins une solution.

Il est dit \textbf{incompatible} s’il n’admet aucune solution.
\end{definition}

\begin{definition}
Un système compatible est dit :
\begin{itemize}
    \item \textbf{déterminé} s’il admet une unique solution ;
    \item \textbf{indéterminé} s’il admet une infinité de solutions.
\end{itemize}
\end{definition}

\begin{exemple}
\[
\begin{cases}
x+y=2,\\
x-y=0
\end{cases}
\]
est déterminé.

\[
\begin{cases}
x+y=2,\\
2x+2y=4
\end{cases}
\]
est indéterminé.

\[
\begin{cases}
x+y=2,\\
2x+2y=5
\end{cases}
\]
est incompatible.
\end{exemple}

\subsection{Interprétation algébrique}

\begin{remarque}
Après réduction, trois types de situations apparaissent souvent :
\begin{itemize}
    \item une ligne du type
    \[
    0=1
    \]
    indique une incompatibilité ;
    \item toutes les inconnues sont déterminées : il y a une solution unique ;
    \item certaines inconnues restent libres : il y a une infinité de solutions.
\end{itemize}
\end{remarque}

\section{Résolution des systèmes triangulaires}

\begin{proposition}
Un système triangulaire dont tous les coefficients diagonaux sont non nuls admet une unique solution.
\end{proposition}

\begin{proof}
Considérons un système triangulaire
\[
\begin{cases}
a_{11}x_1+a_{12}x_2+\cdots+a_{1n}x_n=b_1,\\
a_{22}x_2+\cdots+a_{2n}x_n=b_2,\\
\hspace{2cm}\vdots\\
a_{nn}x_n=b_n,
\end{cases}
\]
avec
\[
a_{11},a_{22},\dots,a_{nn}\neq 0.
\]

La dernière équation donne immédiatement
\[
x_n=\frac{b_n}{a_{nn}}.
\]
En reportant dans l’avant-dernière, on obtient \(x_{n-1}\), puis ainsi de suite jusqu’à \(x_1\). À chaque étape, le coefficient directeur de l’inconnue à déterminer est non nul, donc la valeur obtenue est unique. Ainsi le système admet une unique solution.
\end{proof}

\begin{exemple}
Résolvons
\[
\begin{cases}
x+2y-z=1,\\
3y+4z=5,\\
2z=6.
\end{cases}
\]
La dernière équation donne
\[
z=3.
\]
Puis
\[
3y+12=5 \Longrightarrow 3y=-7 \Longrightarrow y=-\frac73.
\]
Enfin
\[
x+2\left(-\frac73\right)-3=1.
\]
Donc
\[
x-\frac{14}{3}-3=1
\Longrightarrow
x-\frac{23}{3}=1
\Longrightarrow
x=\frac{26}{3}.
\]
\end{exemple}

\section{Systèmes homogènes et structure de l’ensemble des solutions}

\subsection{Stabilité par combinaison linéaire}

\begin{proposition}
L’ensemble des solutions d’un système linéaire homogène est un sous-espace vectoriel de \(\K^n\).
\end{proposition}

\begin{proof}
Considérons le système homogène
\[
\begin{cases}
a_{11}x_1+\cdots+a_{1n}x_n=0,\\
\hspace{2.2cm}\vdots\\
a_{p1}x_1+\cdots+a_{pn}x_n=0.
\end{cases}
\]

Notons \(S\) son ensemble de solutions.

\textbf{1.} Le vecteur nul appartient à \(S\).

\textbf{2.} Soient \(u=(u_1,\dots,u_n)\in S\) et \(v=(v_1,\dots,v_n)\in S\). Alors, pour chaque équation,
\[
a_{i1}u_1+\cdots+a_{in}u_n=0
\quad \text{et} \quad
a_{i1}v_1+\cdots+a_{in}v_n=0.
\]
En additionnant, on obtient que \(u+v\) est encore solution.

\textbf{3.} Soit \(\lambda\in\K\) et \(u\in S\). En multipliant chaque égalité par \(\lambda\), on voit que \(\lambda u\in S\).

Ainsi \(S\) est un sous-espace vectoriel de \(\K^n\).
\end{proof}

\begin{remarque}
C’est un résultat fondamental : les systèmes homogènes conduisent naturellement à des sous-espaces vectoriels.
\end{remarque}

\subsection{Exemple}

\begin{exemple}
Considérons
\[
\begin{cases}
x+y+z=0,\\
x-y=0.
\end{cases}
\]
La seconde équation donne
\[
x=y.
\]
En remplaçant dans la première :
\[
2x+z=0 \Longrightarrow z=-2x.
\]
Donc les solutions sont
\[
(x,x,-2x)=x(1,1,-2).
\]
Ainsi l’ensemble des solutions est
\[
\Vect((1,1,-2)).
\]
C’est une droite vectorielle de \(\R^3\).
\end{exemple}

\section{Systèmes non homogènes et structure de l’ensemble des solutions}

\subsection{Différence de deux solutions}

\begin{proposition}
Soient \(u\) et \(v\) deux solutions d’un même système linéaire non homogène. Alors \(u-v\) est solution du système homogène associé.
\end{proposition}

\begin{proof}
Considérons un système
\[
A x=b.
\]
Si \(u\) et \(v\) sont deux solutions, alors
\[
Au=b
\qquad \text{et} \qquad
Av=b.
\]
En soustrayant,
\[
A(u-v)=0.
\]
Donc \(u-v\) est solution du système homogène associé.
\end{proof}

\subsection{Description de toutes les solutions}

\begin{theoreme}
Si un système linéaire non homogène admet une solution particulière \(u_0\), alors l’ensemble de ses solutions est
\[
u_0+S_h,
\]
où \(S_h\) désigne l’ensemble des solutions du système homogène associé.
\end{theoreme}

\begin{proof}
Soit \(S\) l’ensemble des solutions du système non homogène, et \(S_h\) celui du système homogène associé.

Soit \(u\in S\). Alors, par la proposition précédente,
\[
u-u_0\in S_h,
\]
donc
\[
u=u_0+(u-u_0)\in u_0+S_h.
\]
Ainsi
\[
S\subset u_0+S_h.
\]

Réciproquement, si \(v\in S_h\), alors
\[
A(u_0+v)=Au_0+Av=b+0=b.
\]
Donc \(u_0+v\in S\). Ainsi
\[
u_0+S_h\subset S.
\]

On conclut que
\[
S=u_0+S_h.
\]
\end{proof}

\begin{remarque}
L’ensemble des solutions d’un système non homogène compatible est donc un \textbf{translaté} d’un sous-espace vectoriel.
\end{remarque}

\subsection{Exemple}

\begin{exemple}
Résolvons
\[
\begin{cases}
x+y+z=1,\\
x-y=0.
\end{cases}
\]
On a encore
\[
x=y.
\]
Donc
\[
2x+z=1 \Longrightarrow z=1-2x.
\]
Les solutions sont
\[
(x,x,1-2x).
\]
En posant \(x=t\),
\[
(t,t,1-2t)=(0,0,1)+t(1,1,-2).
\]
Ainsi l’ensemble des solutions est
\[
(0,0,1)+\Vect((1,1,-2)).
\]
\end{exemple}

\section{Systèmes à paramètres}

\subsection{Principe}

\begin{remarque}
Dans de nombreux problèmes, les coefficients ou les seconds membres dépendent d’un paramètre. Il faut alors discuter selon les valeurs de ce paramètre.
\end{remarque}

\begin{methode}
Pour discuter un système à paramètre :
\begin{enumerate}
    \item on applique la méthode du pivot de Gauss comme d’habitude ;
    \item on repère les expressions qu’on divise ou qu’on utilise comme pivots ;
    \item on distingue les cas selon que ces expressions sont nulles ou non.
\end{enumerate}
\end{methode}

\subsection{Exemple}

\begin{exemple}
Discutons selon \(m\in\R\) le système
\[
\begin{cases}
x+y=1,\\
x+my=2.
\end{cases}
\]

On soustrait la première équation à la deuxième :
\[
(m-1)y=1.
\]

\textbf{Cas 1 :} si \(m\neq 1\), alors
\[
y=\frac{1}{m-1},
\]
puis
\[
x=1-y=1-\frac{1}{m-1}=\frac{m-2}{m-1}.
\]
Il y a une unique solution.

\textbf{Cas 2 :} si \(m=1\), le système devient
\[
\begin{cases}
x+y=1,\\
x+y=2,
\end{cases}
\]
ce qui est impossible. Il n’y a aucune solution.
\end{exemple}

\section{Écriture matricielle d’un système}

\subsection{Matrice des coefficients}

\begin{definition}
Au système
\[
\begin{cases}
a_{11}x_1+\cdots+a_{1n}x_n=b_1,\\
\hspace{2.2cm}\vdots\\
a_{p1}x_1+\cdots+a_{pn}x_n=b_p
\end{cases}
\]
on associe :
\begin{itemize}
    \item la matrice des coefficients
    \[
    A=(a_{ij})\in M_{p,n}(\K),
    \]
    \item le vecteur inconnu
    \[
    X=\begin{pmatrix}x_1\\ \vdots\\ x_n\end{pmatrix},
    \]
    \item le vecteur second membre
    \[
    B=\begin{pmatrix}b_1\\ \vdots\\ b_p\end{pmatrix}.
    \]
\end{itemize}
\end{definition}

\begin{definition}
Le système s’écrit alors
\[
AX=B.
\]
\end{definition}

\begin{remarque}
Cette écriture sera fondamentale dans les chapitres suivants.
\end{remarque}

\subsection{Matrice augmentée}

\begin{definition}
La \textbf{matrice augmentée} du système est la matrice obtenue en accolant le second membre à la matrice des coefficients :
\[
\left(
\begin{array}{ccc|c}
a_{11} & \cdots & a_{1n} & b_1\\
\vdots &        & \vdots & \vdots\\
a_{p1} & \cdots & a_{pn} & b_p
\end{array}
\right).
\]
\end{definition}

\begin{remarque}
Les opérations élémentaires se lisent alors comme des opérations sur les lignes de la matrice augmentée.
\end{remarque}

\section{Analyse-synthèse dans les systèmes linéaires}

\begin{remarque}
La méthode d’analyse-synthèse est très naturelle pour les systèmes linéaires :
\begin{itemize}
    \item l’analyse consiste à supposer qu’une solution existe et à dériver sa forme ;
    \item la synthèse consiste à vérifier que les formes trouvées conviennent effectivement.
\end{itemize}
\end{remarque}

\begin{exemple}
Résolvons
\[
\begin{cases}
x+y+z=0,\\
x-y=0.
\end{cases}
\]

\textbf{Analyse :} de \(x-y=0\), on déduit
\[
x=y.
\]
En remplaçant dans la première :
\[
2x+z=0 \Longrightarrow z=-2x.
\]
Donc toute solution éventuelle est de la forme
\[
(x,x,-2x)=t(1,1,-2).
\]

\textbf{Synthèse :} vérifions qu’un vecteur de cette forme est bien solution :
\[
t+t-2t=0
\qquad \text{et} \qquad
t-t=0.
\]
Donc toutes ces valeurs conviennent.

L’ensemble des solutions est donc
\[
\Vect((1,1,-2)).
\]
\end{exemple}

\section{Lien avec les chapitres suivants}

\begin{remarque}
Les systèmes linéaires interviennent immédiatement dans :
\begin{itemize}
    \item le calcul matriciel ;
    \item les espaces vectoriels ;
    \item les familles libres et génératrices ;
    \item la dimension ;
    \item les applications linéaires.
\end{itemize}
\end{remarque}

\begin{exemple}
Pour savoir si une famille \((u_1,\dots,u_p)\) est libre, on résout
\[
\lambda_1u_1+\cdots+\lambda_pu_p=0.
\]
C’est un système linéaire homogène en les \(\lambda_i\).

Pour savoir si elle est génératrice, on cherche à résoudre
\[
x=\lambda_1u_1+\cdots+\lambda_pu_p
\]
pour tout vecteur \(x\). C’est encore un système linéaire.
\end{exemple}

\section{Méthodes pratiques}

\begin{methode}
Pour résoudre un système linéaire :
\begin{enumerate}
    \item écrire proprement le système ;
    \item choisir une stratégie d’élimination ;
    \item effectuer des opérations élémentaires justifiées ;
    \item obtenir un système triangulaire si possible ;
    \item résoudre par remontée ;
    \item vérifier éventuellement la cohérence du résultat.
\end{enumerate}
\end{methode}

\begin{methode}
Pour un système à paramètre :
\begin{enumerate}
    \item ne jamais diviser par une expression sans discuter le cas où elle s’annule ;
    \item distinguer clairement les cas ;
    \item conclure pour chaque cas : aucune solution, une solution unique, ou infinité de solutions.
\end{enumerate}
\end{methode}

\section{Erreurs classiques}

\begin{itemize}
    \item Effectuer une opération sur une équation sans la faire sur le second membre.
    \item Diviser par une quantité qui peut être nulle sans discuter les cas.
    \item Confondre système homogène et système non homogène.
    \item Oublier qu’un système homogène a toujours la solution nulle.
    \item Mal interpréter une ligne du type
    \[
    0=0
    \]
    qui ne donne pas de contrainte nouvelle.
    \item Mal interpréter une ligne du type
    \[
    0=1
    \]
    qui rend le système incompatible.
    \item Conclure trop vite qu’un système est déterminé sans avoir traité toutes les inconnues.
\end{itemize}

\section{Résumé du chapitre}

\begin{itemize}
    \item Un système linéaire est un ensemble d’équations linéaires portant sur les mêmes inconnues.
    \item Les opérations élémentaires conservent l’ensemble des solutions.
    \item La méthode du pivot de Gauss permet de simplifier un système jusqu’à une forme triangulaire.
    \item Un système homogène admet toujours la solution nulle.
    \item L’ensemble des solutions d’un système homogène est un sous-espace vectoriel.
    \item Si un système non homogène est compatible, l’ensemble de ses solutions est un translaté de l’ensemble des solutions du système homogène associé.
    \item Les systèmes à paramètres se discutent selon les valeurs des pivots éventuels.
    \item Un système s’écrit matriciellement sous la forme
    \[
    AX=B.
    \]
\end{itemize}

\section*{Exercices d’application immédiate}

\begin{enumerate}[label=\textbf{Exercice \arabic*.}]
    \item Résoudre dans \(\R^2\) :
    \[
    \begin{cases}
    x+y=5,\\
    2x-y=1.
    \end{cases}
    \]

    \item Résoudre dans \(\R^3\) :
    \[
    \begin{cases}
    x+y+z=2,\\
    x-y=0,\\
    z=1.
    \end{cases}
    \]

    \item Résoudre le système homogène
    \[
    \begin{cases}
    x+y+z=0,\\
    x-y=0.
    \end{cases}
    \]

    \item Décrire l’ensemble des solutions du système
    \[
    \begin{cases}
    x+y+z=1,\\
    x-y=0.
    \end{cases}
    \]

    \item Discuter selon \(m\in\R\) le système
    \[
    \begin{cases}
    x+y=1,\\
    x+my=2.
    \end{cases}
    \]

    \item Écrire sous forme matricielle le système
    \[
    \begin{cases}
    2x-y+z=3,\\
    x+4y-2z=0.
    \end{cases}
    \]

    \item Montrer que l’ensemble des solutions d’un système homogène est un sous-espace vectoriel.

    \item Soient \(u\) et \(v\) deux solutions d’un système non homogène. Montrer que \(u-v\) est solution du système homogène associé.

    \item Résoudre par pivot de Gauss :
    \[
    \begin{cases}
    x+2y-z=1,\\
    2x+y+z=4,\\
    3x+3y=5.
    \end{cases}
    \]

    \item Interpréter géométriquement dans \(\R^2\) les trois systèmes suivants :
    \[
    \begin{cases}
    x+y=1,\\
    x-y=0,
    \end{cases}
    \qquad
    \begin{cases}
    x+y=1,\\
    2x+2y=2,
    \end{cases}
    \qquad
    \begin{cases}
    x+y=1,\\
    2x+2y=3.
    \end{cases}
    \]
\end{enumerate}

\end{document}